\documentclass[11pt]{article}

\usepackage[margin=1in]{geometry}
\usepackage[T1]{fontenc}
\usepackage{lmodern}
\usepackage{amsmath,amssymb}
\usepackage{booktabs}
\usepackage[numbers]{natbib}
\usepackage[hidelinks]{hyperref}
\usepackage{microtype}
\usepackage{enumitem}
\setlist{nosep}
\setlength{\parskip}{0.45em}
\setlength{\parindent}{0pt}

\title{Three Ways the Published Values Can Reveal the Secrets}
\author{A step-by-step cryptanalysis of Chintamani, Paul, and Sa (2024)}
\date{}

\begin{document}
\maketitle

\noindent
\textbf{Goal.} We examine three different ways the secret-masking rule
in the paper can reveal secrets. The order is: (1) one secret is leaked,
(2) no secret is leaked but two public values are available, and
(3) only one public value of the corrected numerical fixture is available.

\medskip
\noindent
\textbf{What these attacks do not do.} None of the attacks below breaks
the elliptic-curve discrete logarithm problem (ECDLP). The ECDLP asks,
in simple terms, whether one can recover a number $a$ from a known curve
point $P$ and the point $aP$. The attacks here do not solve that problem
or recover a hidden scalar from a curve point. Instead, they exploit the
algebraic relationships between the values the scheme publishes, together
with known constraints on the secrets and masks. The paper discusses the
assumed difficulty of the ECDLP in its scheme description
\cite[Sec.~3.1, p.~1459]{chintamani2024}. Our attacks target the published
masking formula, not the underlying ECDLP.

\medskip
\noindent
\textbf{Scope of the results.} The first two attacks follow from the
paper's published-value formula under the conditions stated below.
The third is demonstrated for the paper's small numerical example under
the corrected extension-field encoding stated in Section~4.
It is a demonstrated recovery for that fixture, not a general guarantee
that one public value uniquely determines a secret; other valid inputs
can leave multiple candidates.
These are mathematical demonstrations of weaknesses in the specified
construction; they do not, by themselves, establish whether these
observations have been published previously.

\section{The scheme's published values}

The paper defines a bilinear pairing $e$ and uses the property
\[
    e(P,iQ)=e(P,Q)^i.
\]
Bilinearity is introduced in Section 2. In Section 3.1, the dealer defines
\[
    Q=\sum_j a_{j1}P,
    \qquad
    v_i=e(P,iQ)+K_i,
\]
and publishes the values $v_i$ \cite[Secs.~2 and~3.1, p.~1459]{chintamani2024}.

For convenience, name the common mask
\[
    z=e(P,Q).
\]
Then the public-value formula becomes
\begin{equation}
    \boxed{v_i=z^i+K_i.}
    \label{eq:public-values}
\end{equation}
The secret $K_i$ is an element of the prime field $\mathbb F_p$ (the paper chooses it from the integers $0,1,\ldots,p-1$). The mask $z$ is in the pairing's order-$\ell$ group $\mu_\ell$ in an extension field. The paper discusses this field compatibility in Remark 2 of Section 3.1 \cite{chintamani2024}.

The important point is that all the public values use powers of the same $z$.

\section{One-key recovery}

\subsection{What does the attacker know?}

Suppose the attacker learns the first secret $K_1$. This is the only secret the attacker needs to know at the start. The value $v_1$ is already public, so Equation~\eqref{eq:public-values} gives
\[
    v_1=z+K_1.
\]
Subtract the known secret from the public value:
\begin{equation}
    \boxed{z=v_1-K_1.}
    \label{eq:recover-z-from-k1}
\end{equation}
Thus the common mask is recovered without reconstructing $Q$ or computing a pairing.

\subsection{Why does this reveal the other secrets?}

For every other published value, Equation~\eqref{eq:public-values} says
\[
    v_i=z^i+K_i.
\]
Rearranging gives
\begin{equation}
    \boxed{K_i=v_i-z^i.}
    \label{eq:recover-all}
\end{equation}
The attacker knows $v_i$ because it is public, and now knows $z$ from Equation~\eqref{eq:recover-z-from-k1}. They can therefore calculate every $K_i$.

\textbf{Conclusion.} If $K_1$ leaks, all the other secrets in this construction are exposed.

\subsection{What if a different indexed secret leaks?}

The same conclusion holds for a leaked $K_t$ whenever $t$ has a multiplicative inverse modulo $\ell$. From the public $v_t$,
\[
    z^t=v_t-K_t.
\]
Choose an integer $d$ such that $td=1+a\ell$ for some integer $a$. Since $z\in\mu_\ell$, we have $z^\ell=1$, and therefore
\[
    (v_t-K_t)^d=(z^t)^d=z^{td}
    =z^{1+a\ell}=z(z^\ell)^a=z.
\]
So $z$ can again be recovered. In the paper's example, $\ell=37$ and there are four secrets, indexed $1$ through $4$. Each of those indices has an inverse modulo $37$. For $K_1$, the simpler subtraction in Equation~\eqref{eq:recover-z-from-k1} is enough.

\section{Two-value recovery}

This attack needs no leaked secret. It uses the first two public values, together with the fact that each $K_i$ lies in $\mathbb F_p$.

\subsection{A field fact that removes the secrets}

The field has characteristic $p$. In such a field,
\begin{equation}
    (x+y)^p=x^p+y^p.
    \label{eq:frobenius}
\end{equation}
To see why, expand $(x+y)^p$ using the binomial theorem. Every middle term contains a binomial coefficient $\binom{p}{j}$ for $1\le j\le p-1$. Since $p$ is prime, each such coefficient is divisible by $p$ and is therefore zero in a field of characteristic $p$. Only $x^p$ and $y^p$ remain.

Also, $K^p=K$ for every $K\in\mathbb F_p$: this follows from Fermat's little theorem for nonzero $K$, and is immediate for $K=0$.

\subsection{Use the first public value}

Write the first two public values as
\[
    v_1=z+K_1,
    \qquad v_2=z^2+K_2.
\]
Raise $v_1$ to the $p$th power and subtract $v_1$:
\begin{align*}
    A&=v_1^p-v_1\\
     &=(z+K_1)^p-(z+K_1)\\
     &=z^p+K_1^p-z-K_1\\
     &=z^p-z.
\end{align*}
The secret cancels. Everything on the left is public, so the attacker can calculate $A$.

\subsection{Use the second public value}

Do the same thing to $v_2$:
\begin{align*}
    B&=v_2^p-v_2\\
     &=(z^2+K_2)^p-(z^2+K_2)\\
     &=z^{2p}-z^2\\
     &=(z^p)^2-z^2.
\end{align*}
From $A=z^p-z$, we get $z^p=z+A$. Substituting this into the expression for $B$ gives
\begin{align*}
    B&=(z+A)^2-z^2\\
     &=z^2+2Az+A^2-z^2\\
     &=2Az+A^2.
\end{align*}
Rearrange:
\begin{equation}
    \boxed{z=\frac{B-A^2}{2A}.}
    \label{eq:two-value-z}
\end{equation}
This formula works when $A\ne0$ and $p$ is odd. Both conditions matter: if $A=0$, this particular division is not possible.

\subsection{\texorpdfstring{When is $A$ nonzero?}{When is A nonzero?}}

We found that $A=z^p-z$. Thus $A\ne0$ exactly when $z^p\ne z$. The elements satisfying $x^p=x$ in the extension field are exactly the elements of $\mathbb F_p$.

Here is a short proof. Every element of $\mathbb F_p$ satisfies $x^p=x$, so the polynomial $X^p-X$ already has $p$ roots. A nonzero polynomial of degree $p$ cannot have more than $p$ roots in a field. Consequently it has no other roots in the extension field. Therefore
\[
    \boxed{A\ne0\quad\Longleftrightarrow\quad z\notin\mathbb F_p.}
\]

A useful sufficient condition comes from the pairing group. If $z\ne1$ and $z\in\mu_\ell$ for a prime $\ell$ that does not divide $p-1$, then $z\notin\mathbb F_p$. Otherwise $z$ would lie in the multiplicative group $\mathbb F_p^*$, whose size is $p-1$, and its order $\ell$ would have to divide $p-1$. This is impossible under the stated condition.

The identity case can also be resolved when $\gcd(\ell,p-1)=1$.
If $A=0$, then $z\in\mathbb F_p^*\cap\mu_\ell$. This intersection has
order $\gcd(\ell,p-1)$, so in that case $z=1$ and $K_i=v_i-1$.
The toy two-value script implements this branch. The BLS12-381 dealer
excludes identity masks, and its attacker rejects $A=0$ as outside that
experiment's inputs.

Once $z$ is recovered from Equation~\eqref{eq:two-value-z}, calculate every secret using
\[
    \boxed{K_i=v_i-z^i.}
\]

\textbf{Conclusion.} Under the conditions above, two public values are enough to recover the mask and then all the secrets. No secret needs to leak first.

\section{One-value recovery: the corrected numerical fixture}

This demonstration uses only the first public value $v_1$ and the known structure of the mask. We work through the paper's numerical example under the corrected encoding, where $p=47$, the mask lies in $\mu_{37}$, and $K_1=8$ \cite[Sec.~5, pp.~1462--1463]{chintamani2024}.

\subsection{What is known?}

The displayed public value is
\[
    v_1=37\alpha^2+28\alpha+12,
\]
where the working Sage reproduction represents the extension field using
\[
    \alpha^3+3\alpha+42=0.
\]
A note on the numerical representation appears at the end of this section.

The attacker knows $v_1$, but not $K_1$. They try a candidate $k\in\mathbb F_{47}$ for the secret. This candidate would imply the mask
\[
    w=v_1-k.
\]
For this to be the real mask, it must satisfy
\begin{equation}
    w^{37}=1,
    \label{eq:order-check}
\end{equation}
because the mask belongs to the order-$37$ group $\mu_{37}$.

Checking Equation~\eqref{eq:order-check} for all $47$ possible values of $k$ would already be a finite, direct test for this small example. We can first narrow the candidates using a cubic equation.

\subsection{A necessary condition: the norm}

For an element $w$ in this degree-three extension, define its norm by
\[
    N(w)=w\,w^{47}\,w^{47^2}.
\]
The three factors are the element and its two Frobenius conjugates. Their product lies in $\mathbb F_{47}$; raising the product to the $47$th power simply cycles the three factors.

If $w$ is a valid mask, then $w^{37}=1$. Hence
\begin{align*}
    N(w)&=w^{1+47+47^2}\\
        &=w^{2257}\\
        &=(w^{37})^{61}=1,
\end{align*}
where $2257=1+47+47^2=37\cdot61$. So every valid mask must satisfy $N(w)=1$.

For the candidate $w=v_1-k$, the condition is
\begin{equation}
    N(v_1-k)=(v_1-k)(v_1^{47}-k)(v_1^{47^2}-k)=1.
    \label{eq:norm-condition}
\end{equation}
The substitutions $k^{47}=k$ and $k^{47^2}=k$ are valid because $k\in\mathbb F_{47}$. In the stated field representation, direct field arithmetic gives
\[
\begin{aligned}
    v_1^{47}&=22+27\alpha+42\alpha^2,\\
    v_1^{47^2}&=15+39\alpha+15\alpha^2.
\end{aligned}
\]
Multiplying the three factors in Equation~\eqref{eq:norm-condition}, and reducing all coefficients modulo $47$ using $\alpha^3=-3\alpha-42$, gives
\[
    N(v_1-k)=12-k^3+2k^2-18k.
\]
Set this equal to $1$ and rearrange:
\[
    k^3-2k^2+18k-11=0\pmod{47}.
\]
The cubic factors as
\[
    \boxed{(k-8)(k-12)(k-29)=0\pmod{47}.}
\]
Thus the norm condition leaves three candidates:
\[
    k\in\{8,12,29\}.
\]
This condition is necessary, but it is not sufficient: it tells us the candidate mask has norm $1$, not yet that its 37th power equals $1$.

\subsection{Check which candidate is a valid mask}

Now test the original requirement $w^{37}=1$ for each remaining candidate. Using the same field representation gives
\[
\begin{array}{c|c}
 k & (v_1-k)^{37}\\ \hline
 8  & 1\\
 12 & 42+5\alpha+41\alpha^2\\
 29 & 18+17\alpha+9\alpha^2
\end{array}
\]
The last two results are not $1$: in this field representation, an element equals $1$ only when its constant coefficient is $1$ and its $\alpha$ and $\alpha^2$ coefficients are both zero. Therefore $k=12$ and $k=29$ are rejected, leaving
\[
    \boxed{K_1=8.}
\]
The attacker has recovered the first secret from one public value and the mask's mathematical constraints. This unique result is established for this numerical example. Other valid
public values can leave multiple keys even after the subgroup test, including
within the same field and subgroup. The implementation reports all surviving
candidates rather than assuming uniqueness.

\textbf{Representation note.} Section 5 of the paper gives $X^3+X+4$ as an example of an irreducible cubic, while the working Sage reproduction used for the displayed point and value calculations uses $X^3+3X+42$. The coefficient calculations above are for the latter representation. Coordinates and coefficients tied to one defining polynomial should not be mixed with another without converting the representation \cite[Sec.~5, pp.~1462--1463]{chintamani2024}.

\section{Summary}

\begin{center}
\begin{tabular}{@{}lll@{}}
\toprule
Attack & Attacker starts with & Main step \\
\midrule
One-key recovery & One known secret $K_1$ & $z=v_1-K_1$ \\
Two-value recovery & Public $v_1,v_2$ & Remove secrets with $p$th powers \\
One-value fixture & Public $v_1$ in the corrected fixture & Norm gives 3 candidates; test $w^{37}=1$ \\
\bottomrule
\end{tabular}
\end{center}

All three analyses use the published relationship $v_i=z^i+K_i$. The first is a consequence of one secret being compromised. The second shows how two public values can be enough under a stated field condition. The third uses additional structure in the paper's particular small numerical example.

\clearpage
\section{Current repository implementation}

The main script, \path{chmss_two_level_demo.sage}, constructs and reconstructs
point shares through Levels 1 and 2 and recovers one secret, $K_1=8$.
Its fixed toy fixture uses $Q=7P$. It does not implement the full
three-level example or submitted-share verification.

The attack directories implement the masking layer separately:
\begin{itemize}
    \item \path{attack/one_key/toy_model/}: two toy keys with $Q=7P$;
    the attacker uses the fixed simulated leak \texttt{LEAKED\_K1=8}.
    \item \path{attack/two_values/toy_model/}: four toy keys with $Q=7P$;
    the attacker receives no compromised key.
    \item \path{attack/one_value/}: only the first published value of the
    corrected $Q=13P$ fixture. The attacker enumerates all 47 possible
    keys, tests their candidate masks' norms, then applies $w^{37}=1$.
    It reproduces the candidates $8,12,29$ and the unique survivor $8$
    calculated in Section~4, without reading the expected key.
    This demonstrates recovery for the corrected numerical fixture;
    other valid inputs can leave multiple candidates.
    \item \path{attack/one_key/bls12_381/} and
    \path{attack/two_values/bls12_381/}: four random 256-bit keys embedded
    in the BLS12-381 prime field, using standard G1/G2 generators and a
    reduced Tate pairing. These examples instantiate the same masking
    relation, rather than the complete hierarchical protocol.
\end{itemize}

Each experiment has \texttt{dealer.sage}, \texttt{attacker.sage}, and a
\texttt{public\_values.txt} transcript beside the scripts. Toy one-key and
two-value transcripts use comma-separated power-basis coefficients;
the one-value transcript contains three space-separated coefficients.
The BLS transcripts use JSON. The public files contain masked values,
not private keys or masks. Dealers overwrite their own public transcript.
The BLS one-key dealer also writes a separate \texttt{dealer\_truth.json}
record for verification and prints the first key to model a compromise.

All experiment scripts run directly from the repository root as
\texttt{sage <path>.sage}, with no command-line flags or prior Sage session.
The BLS one-key attacker prompts for the compromised first key. Its verifier
supplies that input automatically. Both \texttt{curve.py} helpers can also
be executed directly with Sage; they define functions rather than run an
attack. Paths and helper imports resolve relative to the script or module.
Commands and full input details are in \path{attack/README.md}.

Only the BLS directories include \texttt{verify\_attack.sage}. Each compares
all four recovered keys with dealer ground truth in two fresh trials.
The two-value verifier uses separate namespaces within one process; the
one-key verifier uses Sage subprocesses and moves the private ground-truth
record away before supplying only $K_1$ to the attacker. It deletes that
record during cleanup. Neither mechanism is an operating-system sandbox.
The toy directories have no separate ground-truth verifier; the one-value
candidate set establishes consistency with the public constraints, and a
singleton set establishes unique recovery under those constraints.

\begin{thebibliography}{1}
\bibitem{chintamani2024}
M.~Chintamani, P.~Paul, and L.~Sa,
``Conjunctive Hierarchical Multi-Secret Sharing Scheme using Elliptic Curves,''
\emph{Indian Journal of Pure and Applied Mathematics}, vol.~55, pp.~1456--1464, 2024.
\href{https://doi.org/10.1007/s13226-023-00450-x}{doi:10.1007/s13226-023-00450-x}.
\end{thebibliography}

\end{document}
